\documentclass[article,a4paper,12pt,oneside,english,brazilian,sumario=tradicional]{abntex2}


\autor{João Pedro Augusto da Silva}
\titulo{Título do artigo}
\local{Brasil}
\data{2026}

% ---
% Pacotes fundamentais
% ---
\usepackage{lmodern}			% Usa a fonte Latin Modern
\usepackage[T1]{fontenc}		% Selecao de codigos de fonte.
\usepackage[utf8]{inputenc}		% Codificacao do documento (conversão automática dos acentos)
\usepackage{indentfirst}		% Indenta o primeiro parágrafo de cada seção.
\usepackage{nomencl} 			% Lista de simbolos
\usepackage{color}				% Controle das cores
\usepackage{graphicx}			% Inclusão de gráficos
\usepackage{microtype} 			% para melhorias de justificação
		
% ---
% Pacotes de citações
% ---
\usepackage[brazilian,hyperpageref]{backref}	 % Paginas com as citações na bibl
\usepackage[alf]{abntex2cite}	% Citações padrão ABNT

% ---
% Pacotes diversos
% ---
\usepackage{lipsum}

\makeindex

\setlrmarginsandblock{3cm}{3cm}{*}
\setulmarginsandblock{3cm}{3cm}{*}
\checkandfixthelayout

% ---
% Espaçamentos entre linhas e parágrafos
% ---

% O tamanho do parágrafo é dado por:
\setlength{\parindent}{1.3cm}

% Controle do espaçamento entre um parágrafo e outro:
\setlength{\parskip}{0.2cm}  % tente também \onelineskip

% Espaçamento simples
\SingleSpacing

% ---
% Documento
% ---
\begin{document}

\selectlanguage{brazilian}
\frenchspacing 

\maketitle

\begin{resumoumacoluna}
    \lipsum[1]

    \vspace{\onelineskip}

    \noindent
    \textbf{Palavras-chave}: artigo. laTeX. neovim.
\end{resumoumacoluna}

% \renewcommand{\resumoname}{Abstract}
% \begin{resumoumacoluna}
%     \begin{otherlanguage*}{english}
%         \lipsum[1]
%
%         \vspace{\onelineskip}
%
%         \noindent
%         \textbf{Keywords}: trees. performance. assyntotic analysis. data structures.
%     \end{otherlanguage*}  
% \end{resumoumacoluna}

\begin{center}
    \smaller
    \textbf{Data de submissão e aprovação}: elemento obrigatório. Indicar dia, mês e ano

    \textbf{Identificação e disponibilidade}: elemento opcional. Pode ser indicado 
    o endereço eletrônico, DOI, suportes e outras informações relativas ao acesso.
\end{center}

\textual

\section{Introdução}

\lipsum[1-2]


\section{Desenvolvimento}

\lipsum[1-5]

Para uma introdução aos algoritmos, ver \cite{exemploLivro}. Aplicações práticas
são discutidas por \citeonline{exemploArtigo} e \citeonline{exemploConferencia},
enquanto uma análise mais profunda pode ser encontrada em \cite{exemploDissertacao}.
Ferramentas de apoio estão documentadas em \citeonline{exemploSite}.



\section{Considerações finais}

\lipsum[1-2]


\postextual

\bibliography{bibliografia/bibliografia}

\end{document}
